%-------------------------------------------------------------------------------
%	SECTION TITLE
%-------------------------------------------------------------------------------
\cvsection{Experience}


%-------------------------------------------------------------------------------
%	CONTENT
%-------------------------------------------------------------------------------
\begin{cventries}

%---------------------------------------------------------
  \cventry
    {Embedded Software Engineer} % Job title
    {Airbus via Critical Software} % Organization
    %{Critical Flytech (formely Critical Software)} % Organization
    {} % Location
    {Jan 2024 - Present} % Date(s)
    {Avionics Platform - Next-generation avionics systems for Airbus commercial aircrafts} % Date(s)
    {
    \begin{cvitems} % Description(s) of tasks/responsibilities
	\item {Engineered Design Assurance Level A and D critical flight software components operating on Real-Time Operating Systems (RTOS), in compliance with DO-178C and DO-330 standards following a V-Model lifecycle. Successfully contributed to major releases, including one achieving Stage of Involvement 3 (SOI-3) certification;}
	\item {Designed and implemented low-level features, such as Aircraft on Ground conditions based on digital signal processing, written in C according to MISRA and compiled with the CompCert formally verified compiler;}
	\item {Managed direct hardware interactions and memory configurations, implementing Direct Memory Access (DMA) transfers and Memory Management Unit (MMU) settings to ensure precise, static memory allocation;}
	\item {Performed formal mathematical verifications (Unit Proof) based on Frama-C/ACSL contracts, alongside automated integration testing using Robot Framework with embedded Remote Procedure Call (eRPC) stubs;}
	\item {Conducted in-depth hardware and memory debugging via TRACE32 to troubleshoot radiation effect statistics and flight data storage issues, while utilizing Wireshark for network packet analysis;}
	\item {Accelerated development workflows by engineering custom Python and Bash scripts for local process automation.}
      \end{cvitems}
    }

%---------------------------------------------------------
  \cventry
    {Data Scientist} % Job title
    {Take The Wind} % Organization
    {} % Location
    {Sep 2022 - Sep 2023} % Date(s)
    {Recommendation system prototype to integrate the main product, Body Interact, to improve users learning rate} % Date(s)
    {
    \begin{cvitems} % Description(s) of tasks/responsibilities
     \item {Performed Requirements Engineering, Risk Management and effort estimation.}
      \item {Developed user clinical competency evaluation algorithm through data and feature engineering, generating 300+ features to expand model training dataset.}
      \item {Created recommendations using cluster-based models (K-Means, Hierachical, DBSCAN) with the aid of healthcare specialists to validate the results.}
      \end{cvitems}
    }

%---------------------------------------------------------
  \cventry
    {Data Scientist Researcher} % Job title
    {Centre for Informatics and Systems of the University of Coimbra} % Organization
    {} % Location
    {Mar 2022 - Sep 2022} % Date(s)
    {CAMELOT - autonomiC plAtform for MachinE Learning using anOnymized daTa for financial fraud detection} % Date(s)
    {
      \begin{cvitems} % Description(s) of tasks/responsibilities
      \item {Benchmarked a privacy-preserving framework for tree-based models (Decision Tree, Random Forest, Gradient-Boosting) with Fully Homomorphic Encryption (BGV, BFV, CKKS via Microsoft SEAL) achieving 88-95\% accuracy on classification and 59-71\% MSE in regression, matching Sklearn baselines across 10 trials;}
        \item {Research paper accepted for presentation at the 27th IEEE Pacific Rim International Symposium on Dependable Computing in November 2022.}
      \end{cvitems}
    }

%---------------------------------------------------------
\end{cventries}
